\section{Uvod}
\label{sec:uvod}

Vektorske baze podataka postale su ključna infrastrukturna komponenta
modernih sistema vještačke inteligencije, omogućavajući semantičku pretragu
nad milijardama visokodimenzionalnih reprezentacija teksta, slika i zvuka.
Kako se ovi sistemi šire iz istraživačkih laboratorija u produkcijska
okruženja, pitanja performansi, memorijske efikasnosti i skalabilnosti
postaju centralna inžinjerska i istraživačka pitanja. Ovaj rad istražuje
kako izbor indeksne strukture i strategija kompresije vektora utiče na
latenciju upita, tačnost pretrage i potrošnju memorije u realnom
eksperimentalnom okruženju.

\subsection{Motivacija}
\label{sec:motivacija}

Primjena velikih jezičkih modela i višemodalnih arhitektura dovela je do
eksplozivnog rasta vektorskih baza podataka u industrijskim sistemima.
Embeddingsi koji se generišu modelima poput BERT-a, CLIP-a i srodnih
arhitektura tipično su dimenzionalnosti od $768$ do $2048$ i pohranjuju se
kao \foreignlanguage{english}{float32} vrijednosti.
Pri pohrani milijardi ovakvih vektora memorijski zahtjevi dosežu i
nekoliko terabajta, što daleko premašuje kapacitet standardnog
hardvera~\cite{raja2026memory}.

Osim memorijskog pritiska, visokodimenzionalni prostori uvode i
računarski izazov poznat kao \emph{prokletstvo dimenzionalnosti}: iscrpna
pretraga nad svim vektorima skalira linearno s veličinom skupa podataka
i postaje nepraktična pri web-razmjernim kolekcijama~\cite{ma2026comprehensive}.
Ovaj problem motiviše razvoj aproksimativnih metoda pretrage koje,
uz kontrolirani gubitak tačnosti, omogućavaju odgovore u milisekundama.
Uprkos dostupnosti naprednih indeksnih struktura i tehnika kompresije,
njihova komparativna analiza u standardizovanim uvjetima ostaje
nedovoljno istražena, posebno u kontekstu kompromisa između latencije,
tačnosti i potrošnje memorije na istom skupu podataka.

\subsection{Indeksi i kvantizacija u vektorskim bazama podataka}
\label{sec:pregled_tehnika}

Savremeni sistemi za vektorsku pretragu oslanjaju se na dvije klase
tehnika: indeksne strukture i metode kompresije vektora.

Indeksne strukture određuju kako su vektori organizovani radi efikasne
pretrage. \foreignlanguage{english}{\emph{Flat}} indeks ne vrši nikakvu
aproksimaciju, svaki upit poredi se s cijelim skupom podataka, što
garantuje tačne rezultate ali skalira s $\mathcal{O}(n)$ složenošću.
\foreignlanguage{english}{\emph{Hierarchical Navigable Small World}}
(HNSW) gradi višeslojni grafovski indeks koji omogućava logaritamsku
složenost pretrage~\cite{malkov2018efficient}.
\foreignlanguage{english}{\emph{Inverted File Index}} (IVF) dijeli
prostor na klastere i pretragu ograničava na podskup
klastera~\cite{jegou2011product}.

Metode kompresije smanjuju memorijski otisak indeksa kodiranjem vektora
u kompaktnije reprezentacije. Bez kompresije (\foreignlanguage{english}{Flat})
svaki vektor zauzima punu \foreignlanguage{english}{float32} preciznost.
\foreignlanguage{english}{\emph{Binary Quantization}} svodi svaku
dimenziju na jedan bit, postižući kompresiju $32\times$.
\foreignlanguage{english}{\emph{Scalar Quantization}} mapira svaku
dimenziju na $b$-bitni cijeli broj, tipično $b=8$, smanjujući pohranu
četverostruko.
\foreignlanguage{english}{\emph{Product Quantization}} dijeli vektore na
pod-prostore koji se kvantiziraju nezavisno, postižući veće kompresijske
omjere uz bolju rekonstrukcijsku tačnost od skalarne
kvantizacije~\cite{raja2026memory}.

\subsection{Ciljevi rada}
\label{sec:ciljevi}

Cilj ovog rada je empirijska evaluacija i komparativna analiza
indeksnih struktura i tehnika kvantizacije za vektorsku pretragu na
\foreignlanguage{english}{MS MARCO} skupu podataka~\cite{msmarco2016},
koji predstavlja standardni benchmark za zadatke pretrage informacija.
Eksperimenti su dizajnirani da kvantificiraju kompromise između
latencije upita, tačnosti pretrage i potrošnje memorije pri različitim
kombinacijama indeksa i kompresije.

\noindent U skladu s tim, postavljaju se sljedeće istraživačke hipoteze:

\begin{description}
    \item[H1:] HNSW indeks postiže statistički značajno nižu latenciju
    upita u poređenju s IVF indeksom na \foreignlanguage{english}{MS MARCO}
    skupu podataka, uz zadržavanje \foreignlanguage{english}{recall@10}
    vrijednosti iznad $90\%$.

    \item[H2:] Primjena \foreignlanguage{english}{Product Quantization}
    tehnike smanjuje memorijski
    \foreignlanguage{english}{\emph{footprint}} vektorskog indeksa za
    najmanje $4\times$ u poređenju s nekomprimiranim
    \foreignlanguage{english}{\emph{flat}} indeksom, uz smanjenje
    \foreignlanguage{english}{recall@10} manje od $5$ procentnih poena.

    \item[H3:] Kombinacija IVF i \foreignlanguage{english}{Product
    Quantization} (IVF-PQ) pruža bolji
    \foreignlanguage{english}{\emph{recall-latency trade-off}} u poređenju
    sa samostalnim HNSW pristupom kada je memorijski budžet ograničen
    na ispod $2$\,GB.
\end{description}
